\documentclass[11pt,a4paper]{article}
\usepackage[T1]{fontenc}
\usepackage{lmodern}
\usepackage{graphicx}
\usepackage{amsmath}
\usepackage{booktabs}
\usepackage{hyperref}
\usepackage{geometry}
\usepackage{caption}
\usepackage{float}
\usepackage{placeins}
\usepackage{parskip}
\geometry{margin=1in}

\hypersetup{
  colorlinks=true,
  linkcolor=black,
  urlcolor=blue,
  pdfauthor={Team Dominators},
  pdftitle={The Sentinel Shield - Near-Earth Comet (NEC) Classification}
}

\setlength{\parskip}{4pt}
\setlength{\parindent}{0pt}
\renewcommand{\baselinestretch}{1.05}

\title{\textbf{\Large The Sentinel Shield - Near-Earth Comet (NEC) Classification}}
\author{
\textbf{\MakeUppercase{TEAM DOMINATORS}}\\
\textbf{Aryan Sisodiya (B25126)}, \textbf{Daksh Rathi (B25132)}, \textbf{Priyanshu Baranwal (B25165)}\\
SpaceCode Hackathon 2026 — STAC IIT Mandi
}
\date{}

\begin{document}
\maketitle

\begin{center}
{\large\textbf{Abstract}}
\end{center}
We present a machine learning system for automatic classification of Near-Earth Objects (NEOs) as
Potentially Hazardous Objects (PHOs) using orbital and physical parameters from NASA’s NEOWS catalog
(1950--2025). The model learns the physical hazard boundary defined by Minimum Orbit Intersection
Distance (MOID) and Absolute Magnitude, while incorporating orbital dynamics such as eccentricity,
inclination and relative velocity. A tree-based XGBoost classifier trained with class imbalance handling
achieves \textbf{ROC-AUC = 0.999999} and \textbf{Recall = 0.9996}. The trained system is deployed as an
interactive web application for real-time planetary defense risk assessment.

\section*{Executive Summary}
\textbf{What we did:} Built a supervised ML model to classify Near-Earth Objects as hazardous or safe.\\\\  
\textbf{Performance:} ROC-AUC = 0.999999, Recall = 0.9996.\\\\  
\textbf{Physics alignment:} Model recovers NASA’s MOID--Magnitude PHO rule.\\\\  
\textbf{Deployment:} Live web app for real-time asteroid hazard prediction.

\section{Problem}
Near-Earth Objects such as asteroids and comets can pose a serious threat when their orbits bring them
close to Earth and their size is sufficiently large. NASA defines a Potentially Hazardous Object (PHO)
as one with Minimum Orbit Intersection Distance (MOID) not exceeding 0.05 AU and Absolute Magnitude
$H$ not exceeding 22. PS-0.3 requires building a reliable and reproducible ML system that predicts this
hazard status from orbital and physical parameters.

\section{Dataset}
We use NASA’s Near-Earth Object Web Service (NEOWS) dataset:
\href{https://drive.google.com/file/d/1BN6ro6Qtfd4j7tUtlrgVzZ_Ts3axj6v0/view}{nasa\_neows\_1950\_2025.csv}.
It contains 375,659 objects with physical properties, orbital elements, proximity metrics and
dynamical indicators spanning 1950 to 2025.

\section{Method}
We remove identifiers and high-cardinality name fields, apply median imputation for missing values,
log-transform skewed physical features such as diameter and velocity, and standardize all features.
Because hazardous objects are statistically rare, SMOTE oversampling is used. An \textbf{XGBoost}
classifier is trained to maximize recall and ROC-AUC.

\section{Compute}
Training was performed on an \textbf{NVIDIA RTX 3050 (6GB)} using Python, pandas, scikit-learn and
XGBoost. Deployment is handled through Streamlit Cloud.

\section{Results}
The classifier separates hazardous and non-hazardous objects with near-perfect accuracy.
A sharp probability transition is observed when MOID crosses the physical threshold of 0.05 AU,
indicating that the learned decision boundary aligns with orbital mechanics.

\section{Deployment}
The trained model is deployed as a live interactive web application where users can test real
asteroids and modify orbital parameters to observe hazard probabilities.

Live application:
\href{https://spacecode-ps3.streamlit.app/}{https://spacecode-ps3.streamlit.app/}

\section{Conclusion}
We built a complete, reproducible and deployed planetary defense classifier that identifies hazardous
Near-Earth Objects from real astronomical data. The Sentinel Shield system demonstrates how machine
learning and orbital mechanics can be combined into a working risk-assessment interface.

\section*{Code and References}
\textbf{Project Repository (PS-3):}
\href{https://github.com/InfinityxR9/SpaceCode}{SpaceCode} (folder \texttt{ps3/})

\textbf{Dataset:}
\href{https://drive.google.com/file/d/1BN6ro6Qtfd4j7tUtlrgVzZ_Ts3axj6v0/view}{NASA NEOWS Dataset}

\textbf{Deployment:}
\href{https://spacecode-ps3.streamlit.app/}{Sentinel Shield Live App}

\textbf{Team GitHub Profiles:}
\begin{itemize}
\item Aryan Sisodiya — \href{https://github.com/InfinityxR9}{InfinityxR9}
\item Daksh Rathi — \href{https://github.com/dakshrathi-india}{dakshrathi-india}
\item Priyanshu Baranwal — \href{https://github.com/priyanshubarnwal122-ai}{priyanshubarnwal122-ai}
\end{itemize}

\end{document}
